\section{模型的评价与推广}
\label{sec:7}

\subsection{模型的优点}
\label{sec:7.1}

\begin{enumerate}
\item \textbf{唯一质量口径与清晰证据分层。} 文档、领域、桥接标签与配比加权质量统一使用 CRITIC 加权 Huber 主分；A1 拟合、A2/A3 扩展、A6/A7 留出、A8/A9 标定、A10/A11 独立检验、A12–A15 估算分层报告。
\item \textbf{训练与检验边界可追溯。} 问题二只用 B6 估计质量参数、选模型，B7 新增 90 点作留出；配比通道参数只由 1M 训练行确定，A6–A11 仅用于评估。
\item \textbf{形式与阈值可检验。} 六来源阈值取 A1 同域 P95，桥接以分层五折比较 Ridge 与 GBDT；结构性转移由质量约束活跃状态定义并区分局部根与全局切换；主式显式指定并同时报告交叉验证更优的线性对照。
\item \textbf{解析与数值互验。} 替代条件、$N^*$ 与算力等价倍数、降维闭式与 $\Phi$ 判据、S 形桥接与 Shapley 分解均有闭式或可核验形式，并与数值结果逐项对照。
\item \textbf{职责分离、避免循环论证。} 先固定主分再报告分歧诊断与旧法对照，阈值敏感性与折扣对照分均不回写主评分，筛查预算不被当作真实冲突率。
\end{enumerate}

\subsection{模型的局限与改进方向}
\label{sec:7.2}

本研究的证据边界集中说明如下，凡正文结论涉及这些条件时均以此为准。

\begin{enumerate}
\item \textbf{无独立质量真值。} 问题一以方法一致性、扰动稳健性与下游敏感性刻画质量口径，不给出质量准确率；17 域中 11 域质量依赖文本桥接（五折 $R^2=0.4106$），区间不含模型选择与跨域偏差，ubuntu\_irc（16 篇）、philpapers（63 篇）低置信使用。
\item \textbf{质量参数来自半合成数据。} $\kappa_N,\kappa_D,k_0$ 只由 B6（半合成，最大约 12B）估计，B7 留出为 $Q$ 方向插值；主式为指定幂次，留出 RMSE（$0.0531$）高于线性对照（$0.0416$）与噪声代理（$0.0458$），$\kappa_D$ 较弱；高预算端地板份额（约 $33\%$）与「提质远比扩参数划算」含地板外推成分。
\item \textbf{跨尺度外推幅度不足。} 1B 检验显示尺度修正把绝对误差降低约 $74\%$，但平均预测仍高于观测，配比效应幅度在 1B 缩至约六成；两档规模下单参数缩放（对数线性或幂律）刻画不出曲率，曲率参数欠定。
\item \textbf{问题三的数据边界。} Q3-9 已闭合固定配比与重选配比两类抽样；质量参数仍来自半合成 B6，配比通道为探索性，$\log_{10}C\gtrsim23.5$ 的最优规模配置为外推。
\item \textbf{问题四口径尚未由重跑固定。} 主面板未执行严格开放权重纳入与许可证审计，C6 桥接迁移误差未单独进入预测区间，预测以历史锚点 $T_0=2025.25$ 为起点且原面板含前视数据，主链未在路径修复后重跑，补充独立实现尚未在本仓库运行（\texttt{DEFERRED\_Q3}）。
\item \textbf{问题四方法分歧与占比口径。} C6 的 High 层仅 7 点且全在底噪区，上升段依赖 Medium 层；分解三口径技术占比极差达 $32.81$ pp，是最大不确定性来源；占比分母为模型解释增量而非观测增量，观测与模型增量不完全闭合。
\item \textbf{预测不含离散突破。} 条件预测不含新架构等非平滑变化；个别低算力情景区间下界低于锚点，与累计前沿不下降定义存在张力，需先修复单调性。
\end{enumerate}

\textbf{改进方向。}
\begin{enumerate}
\item 为 11 个 inferred 域增补质量标签或模型评分信号，重新校准并检验桥接误差；补齐 P1 逐次重抽样接口，使两类配比不确定性可分别传播到问题三。
\item 取得真实跨质量、跨规模训练日志后替换 B6 重估 $k_0$ 与 $\kappa_N$，并允许配比强度随 $\log N$ 变化；在获得真正跨尺度配比试验前对 H$_{\rm red}$、H$_{\rm tot}$ 各算一遍。
\item 问题四重跑时执行明确的开放权重纳入规则（报告纳入前后样本量与敏感性），把建模、校准与回测样本全部截在各自起点之前，将 C6 桥接的 bootstrap 与残差作为独立不确定性通道，并统一占比分母与前沿单调性。
\item 取得更多同验证集的 Loss–Benchmark 对，缩小结构口径与线性口径的技术占比分歧。
\end{enumerate}

\subsection{模型的推广}
\label{sec:7.3}

「经典基线＋质量修正＋配比乘子」的标度律结构可推广到其他数据受限的资源配置问题（去重、合成数据、多模态对齐均可写成可约损失上的乘子与不随规模消失的地板）；「约束取等后解析消元＋活跃约束状态分类」的降维方法适用于含多种数据处理成本的预算优化；「分层 S 形桥接＋连续时间项＋Shapley 分解」框架适用于其他存在能力饱和的技术演进分析。推广时保持同一纪律：明确数据证据等级，把半合成、估算与外推数据标注为假设支撑而非观测验证。
